%==============================================================================
\chapter{Diseño experimental}
\label{cap:diseno}
%==============================================================================
% Objetivos específicos 1 y 4.

Para comprobar si el reparto del tiempo entre configuraciones y el uso de las
características de cada instancia mejoran la resolución del JSSP frente a
ejecutar siempre un único \textit{solver}, se diseñó una evaluación controlada
sobre las instancias de referencia. Antes de pasar al análisis cuantitativo del
Capítulo~\ref{cap:resultados}, aquí se precisan las preguntas experimentales,
el esquema de validación cruzada, los métodos comparados y las métricas
empleadas.

%------------------------------------------------------------------------------
\section{Preguntas de evaluación}
\label{sec:preguntas}
%------------------------------------------------------------------------------

La evaluación se organizó en torno a cinco preguntas, derivadas de los
objetivos específicos:

\begin{enumerate}[label=\textbf{P\arabic*.}, leftmargin=*]
  \item \textbf{Efecto de las cotas.} ¿Cuánto mejora el desempeño de los
    \textit{solvers} al pasar del modelo original al modelo con cotas?
  \item \textbf{Complementariedad.} ¿Las configuraciones del portafolio son
    complementarias? ¿En qué tipo de instancias es mejor cada una y cuánto
    margen deja eso para un portafolio?
  \item \textbf{\textit{Winner Takes All} frente a \textit{schedule}
    secuencial} (objetivo específico~1). ¿Repartir el tiempo entre varias
    configuraciones da mejores resultados que elegir una sola?
  \item \textbf{Aporte de las características y del aprendizaje automático}
    (objetivo específico~3). ¿Las estrategias que usan las características
    de la instancia superan a las que no las usan, como el mejor
    \textit{solver} individual, la regla por tamaño o el \textit{schedule}
    estático?
  \item \textbf{Comparación con el estado del arte} (objetivo
    específico~4). ¿Cómo se sitúa el portafolio frente a CP-SAT, un
    \textit{solver} de referencia que no forma parte de él?
\end{enumerate}

%------------------------------------------------------------------------------
\section{Instancias y esquema de validación}
\label{sec:validacion}
%------------------------------------------------------------------------------

La evaluación usa las 236 instancias en las que $lb < ub$. Las seis restantes
quedan resueltas por las cotas (Sección~\ref{sec:cotas_instancias}), puesto que
todos los métodos obtienen en ellas el óptimo sin ejecutar ningún
\textit{solver}, así que incluirlas no cambiaría la comparación.

Como el conjunto de instancias es pequeño, se usó validación cruzada dejando
una instancia fuera (\textit{leave-one-out}). Para cada una de las 236
instancias se toman las otras 235 como entrenamiento, se calcula con ellas
todo lo que el método necesita aprender (la imputación y la normalización de
las características, los vecinos más cercanos, el umbral de la regla por
tamaño, el \textit{schedule} estático y los modelos de aprendizaje automático)
y se construye el \textit{schedule} de la instancia excluida, que se evalúa con
sus resultados reales. Así, ninguna decisión sobre una instancia usa
información de esa misma instancia, y las comparaciones entre métodos son
pareadas, porque todos se evalúan sobre las mismas 236 instancias.

\subsection{Evaluación de \textit{schedules} por simulación}

Cada \textit{schedule} se evaluó directamente sobre las trayectorias
registradas en la campaña con cotas, mediante la
ecuación~\eqref{eq:resultado_schedule}. Para una configuración $a$ con un
tiempo asignado $\tau$, su resultado es el mejor \textit{makespan} que $a$
había encontrado en el instante $\tau$ de su ejecución de 300 segundos, o la
$ub$ si todavía no había encontrado ninguno. Si $a$ probó el óptimo antes de
$\tau$, el \textit{schedule} termina en ese momento. Esta simulación es exacta
porque las configuraciones de un \textit{schedule} se ejecutan de forma
independiente, por lo cual la ejecución de $a$ durante $\tau$ segundos dentro
del \textit{schedule} coincide con el tramo inicial de su ejecución de 300
segundos.

La simulación permite comparar muchas estrategias y variantes con el mismo
costo computacional que una sola campaña, y garantiza que las diferencias
entre métodos se deban a sus decisiones y no a variaciones entre ejecuciones.

%------------------------------------------------------------------------------
\section{Enfoques comparados y referencias}
\label{sec:enfoques}
%------------------------------------------------------------------------------

En la Tabla~\ref{tab:metodos} se clasifican los métodos evaluados, cuya
implementación fue descrita en el Capítulo~\ref{cap:portafolio}. Se agrupan en
referencias, estrategias que no usan características, estrategias de selección
de una configuración (\textit{Winner Takes All}) y estrategias que construyen
un \textit{schedule}.

WTA $k$-NN y SUNNY se evaluaron con los tres conjuntos de características de
la Sección~\ref{sec:features}, formados por las de \texttt{mzn2feat}, las
propias y la unión de ambas. El \textit{schedule} con ML usa la unión. El
número de vecinos es $k = 15$, y la sensibilidad a este valor se evaluó con
$k \in \{3, 5, 7, 10, 15, 20, 30\}$.

El VBS no es un método aplicable, porque requiere conocer de antemano el
resultado de cada configuración, de modo que representa la cota superior de lo
que puede lograr cualquier estrategia con las mismas configuraciones. CP-SAT,
en cambio, es un \textit{solver} real, pero no forma parte del portafolio
(Sección~\ref{sec:composicion}) y su papel consiste en situar los resultados
respecto al estado del arte.

\begin{table}[H]
  \centering
  \small
  \singlespacing
  \caption{Métodos comparados.}
  \label{tab:metodos}
  \begin{tabularx}{\textwidth}{l >{\raggedright\arraybackslash}X c c}
    \toprule
    \textbf{Método} & \textbf{Descripción} & \textbf{Usa caract.} & \textbf{Varias config.} \\
    \midrule
    \multicolumn{4}{l}{\textit{Referencias}} \\
    Configuraciones individuales & Chuffed, Gecode y Chuffed V2 durante 300\,s & No & No \\
    SBS & La configuración con menor \textit{gap} medio en el entrenamiento & No & No \\
    VBS & La mejor configuración en cada instancia (oráculo) & -- & No \\
    CP-SAT & \textit{Solver} de referencia externo, fuera del portafolio & No & No \\
    \midrule
    \multicolumn{4}{l}{\textit{Sin características}} \\
    Regla por tamaño & Una configuración para instancias pequeñas y otra para grandes & Solo $n \cdot m$ & No \\
    \textit{Schedule} estático & El mismo \textit{schedule} para todas las instancias & No & Sí \\
    \midrule
    \multicolumn{4}{l}{\textit{Winner Takes All}} \\
    WTA $k$-NN & La configuración con menor \textit{gap} medio en los $k$ vecinos & Sí & No \\
    \midrule
    \multicolumn{4}{l}{\textit{Schedule secuencial}} \\
    SUNNY & \textit{Schedule} a partir de los $k$ vecinos más cercanos & Sí & Sí \\
    \textit{Schedule} con ML & \textit{Schedule} a partir de tres modelos de predicción & Sí & Sí \\
    \bottomrule
  \end{tabularx}
\end{table}

%------------------------------------------------------------------------------
\section{Métricas}
\label{sec:metricas}
%------------------------------------------------------------------------------

Se usaron las métricas de la Sección~\ref{sec:metricas_portafolios}, con
$T = 300$ segundos. La principal es el \textbf{\textit{gap} medio}
(ecuación~\eqref{eq:gap}), porque solo 36 de las 236 instancias se pueden
probar óptimas dentro del límite con alguna configuración y, en la gran
mayoría, lo relevante es la calidad de la mejor solución encontrada. Se
reportan también los \textbf{óptimos probados} y el \textbf{PAR10}
\citep{lindauer2019}, con
$10 \cdot T = 3000$ segundos para las instancias en que no se prueba el óptimo.
En un \textit{schedule}, el tiempo incluye el de las configuraciones
ejecutadas antes de la que lo prueba. La \textbf{brecha cerrada} es la
fracción del margen entre el SBS y el VBS que el método logra cerrar
\citep{kerschke2019, lindauer2019},
\begin{equation}
  \mathit{brecha} = 100 \cdot
  \frac{\overline{\mathit{gap}}_{\mathrm{SBS}} - \overline{\mathit{gap}}}
       {\overline{\mathit{gap}}_{\mathrm{SBS}} - \overline{\mathit{gap}}_{\mathrm{VBS}}},
  \label{eq:brecha}
\end{equation}
que vale 0\,\% para el SBS, 100\,\% para el VBS y es negativa si el método es
peor que usar siempre la mejor configuración individual. Por último, el número
de instancias en que el método obtiene un \textit{makespan} menor o mayor que
el SBS distingue a un método que mejora mucho en pocas instancias de uno que
mejora poco en muchas.

Para comparar dos métodos se usan sus diferencias de \textit{gap} instancia
por instancia. Se reportan la diferencia media con un intervalo de confianza
del 95\,\% obtenido por \textit{bootstrap} (10\,000 remuestreos) y el
$p$-valor de la prueba de rangos con signo de Wilcoxon sobre las instancias
en que los resultados difieren. El análisis descriptivo de la campaña incluye
asimismo el \textit{gap} medio de cada configuración en función del tiempo (10,
30, 60, 120 y 300 segundos), que muestra su comportamiento \textit{anytime}.

%------------------------------------------------------------------------------
\section{Protocolo de ejecución}
\label{sec:protocolo}
%------------------------------------------------------------------------------

Los datos provienen de la campaña con cotas de la
Sección~\ref{sec:ejecucion_solvers}, en la cual cada configuración se ejecutó
una vez sobre cada instancia, con un límite de 300 segundos y un solo núcleo.
Las ejecuciones se hicieron de dos en dos, cada una fijada a un núcleo físico
distinto del procesador. Las versiones del software y las características del
equipo se detallan en el Anexo~\ref{anx:configuracion}.

Los métodos de selección y los modelos de aprendizaje se implementaron en
Python con scikit-learn \citep{pedregosa2011}. Los bosques aleatorios usan una
semilla fija, y los empates en los vecinos más cercanos y en la elección de
configuraciones se resuelven con un orden fijo, de modo que todos los
resultados son deterministas.

%------------------------------------------------------------------------------
\section{Reproducibilidad}
\label{sec:reproducibilidad}
%------------------------------------------------------------------------------

Todo el proceso puede repetirse con el repositorio que acompaña este
documento, cuyos programas siguen el orden del trabajo. Los scripts \texttt{01}
a \texttt{04} convierten las instancias y calculan las cotas y las
características (Capítulo~\ref{cap:caracterizacion}), \texttt{05} ejecuta las
campañas, \texttt{06} construye la base de conocimiento y el análisis
descriptivo, \texttt{07} y \texttt{08} implementan las estrategias de selección
y los modelos de aprendizaje con su validación, y \texttt{09} genera las tablas
y figuras del Capítulo~\ref{cap:resultados}. Como los resultados de las
campañas se incluyen en el repositorio, los programas \texttt{06} a \texttt{09}
pueden ejecutarse sin repetirlas, y la mayor parte de su tiempo de cómputo
corresponde al entrenamiento de los modelos en cada partición de la
validación. La guía completa está en el Anexo~\ref{anx:reproduccion}.
